% English translation of chapter7.tex lines 1130--1279.
% Source: Methods of Algebra, Volume 2, commit
% 9a5803ff77dd3257484cb177f851a73770a59dd3 (CC BY 4.0).
% stable-unit-id: o014.aljabr2.chapter7.coalgebra-to-hopf-b
% Coverage: the remainder of Section 7.5, from after the proof of the
% monoidal structure on a bialgebra's module category to line 1279.

% segment-id: o014.aljabr2.chapter7.coalgebra-to-hopf-b.p001
For a bialgebra $A$, one can still ask whether the monoidal category
$A\dcate{Mod}$ has a braiding and whether that braiding is symmetric. A
natural first idea is to use the braiding $c(M_1,M_2)$ of $\mathcal{V}$, but
in general this is not a morphism in $A\dcate{Mod}$. Even in the special case
$\mathcal{V}=\Bbbk\dcate{Mod}$, the answer already involves the concept of a
\emph{quasitriangular bialgebra}, which cannot be explained briefly. We give
only the following example, which is so simple as to be somewhat misleading.

% segment-id: o014.aljabr2.chapter7.coalgebra-to-hopf-b.pr001
\begin{proposition}\label{prop:bialgebra-mod-symmetry}
	Let $\mathcal{V}$ be a symmetric monoidal category.
% segment-id: o014.aljabr2.chapter7.coalgebra-to-hopf-b.l001
	\begin{enumerate}[(i)]
% segment-id: o014.aljabr2.chapter7.coalgebra-to-hopf-b.i001
		\item If $A$ is a cocommutative bialgebra over $\mathcal{V}$, then the
		monoidal categories $A\dcate{Mod}$ and $\cated{Mod}A$ become symmetric
		monoidal categories with respect to
		$c(M_1,M_2):M_1\otimes M_2\rightiso M_2\otimes M_1$.

% segment-id: o014.aljabr2.chapter7.coalgebra-to-hopf-b.i002
		\item Dually, if $A$ is a commutative bialgebra over $\mathcal{V}$,
		then $A\dcate{Comod}$ and $\cated{Comod}A$ are symmetric monoidal
		categories.
	\end{enumerate}
\end{proposition}

% segment-id: o014.aljabr2.chapter7.coalgebra-to-hopf-b.pf001
\begin{proof}
	It is enough to treat the cocommutative case. Using the preceding notation,
	the assertion that $c(M_1,M_2)$ is a morphism in $A\dcate{Mod}$ is
	equivalent to proving that the outer boundary of the following diagram
	commutes:
% segment-id: o014.aljabr2.chapter7.coalgebra-to-hopf-b.g001
	\[\begin{tikzcd}[row sep=large]
		A \otimes M_1 \otimes M_2 \arrow[r, "{\Delta \otimes \identity \otimes \identity}" inner sep=0.6em] \arrow[d, "{\identity \otimes c(M_1, M_2)}" description] & A \otimes A \otimes M_1 \otimes M_2 \arrow[r, "{\identity \otimes c(A, M_1) \otimes \identity}" inner sep=0.6em] \arrow[d, "{c(A, A) \otimes c(M_1, M_2)}" description] & A \otimes M_1 \otimes A \otimes M_2 \arrow[r, "{\mu_1 \otimes \mu_2}" inner sep=0.6em] \arrow[d, "{c(A \otimes M_1, A \otimes M_2)}" description] & M_1 \otimes M_2 \arrow[d, "{c(M_1, M_2)}" description] \\
		A \otimes M_2 \otimes M_1 \arrow[r, "{\Delta \otimes \identity \otimes \identity}"' inner sep=0.6em] & A \otimes A \otimes M_2 \otimes M_2 \arrow[r, "{\identity \otimes c(A, M_2) \otimes \identity}"' inner sep=0.6em] & A \otimes M_2 \otimes A \otimes M_1 \arrow[r, "{\mu_2 \otimes \mu_1}"' inner sep=0.6em] & M_2 \otimes M_1 .
	\end{tikzcd}\]

	The square on the left commutes because $A$ is cocommutative, while the
	square on the right commutes by functoriality of $c$. Drawing the braid
	diagram shows that
% segment-id: o014.aljabr2.chapter7.coalgebra-to-hopf-b.g002
	\begin{multline*}
		c(A \otimes M_1, A \otimes M_2) = \\
		(\identity_A \otimes c(A, M_2) \otimes \identity_{M_1}) (c(A, A) \otimes c(M_1, M_2)) (\identity_A \otimes c(A, M_1) \otimes \identity_{M_2}),
	\end{multline*}
	and since $\mathcal{V}$ is a symmetric monoidal category, the middle
	square also commutes. All other properties of the symmetric monoidal
	structure on $A\dcate{Mod}$ reduce to verifications in $\mathcal{V}$.
\end{proof}

% segment-id: o014.aljabr2.chapter7.coalgebra-to-hopf-b.p002
We now introduce the concept of a Hopf algebra. Let $\mathcal{V}$ be a
braided monoidal category.

% segment-id: o014.aljabr2.chapter7.coalgebra-to-hopf-b.d001
\begin{definition}\label{def:Hopf-algebra}
% segment-id: o014.aljabr2.chapter7.coalgebra-to-hopf-b.x001
	\index{Hopf algebra}
% segment-id: o014.aljabr2.chapter7.coalgebra-to-hopf-b.x002
	\index{antipode}
	Let the data $(A,\mu,\eta,\Delta,\epsilon)$ form a bialgebra $A$ over
	$\mathcal{V}$. If an isomorphism $S:A\rightiso A$ in $\mathcal{V}$ makes
	the following diagram commute,
% segment-id: o014.aljabr2.chapter7.coalgebra-to-hopf-b.g003
	\[\begin{tikzcd}
		A \otimes A \arrow[d, "\mu"'] & A \otimes A \arrow[r, "{S \otimes \identity_A}"] \arrow[l, "{\identity_A \otimes S}"'] & A \otimes A \arrow[d, "\mu"] \\
		A & A \arrow[u, "\Delta"] \arrow[r, "\eta\epsilon"'] \arrow[l, "\eta\epsilon"] & A ,
	\end{tikzcd}\]
	that is,
% segment-id: o014.aljabr2.chapter7.coalgebra-to-hopf-b.q001
	\[ \mu (\identity_A \otimes S) \Delta = \eta\epsilon = \mu (S \otimes \identity_A) \Delta, \]
	then $S$ is called an \emph{antipode} of $A$. A bialgebra equipped with an
	antipode is called a \emph{Hopf algebra}. Morphisms between Hopf algebras
	are defined to be morphisms between their underlying bialgebras.
\end{definition}

% segment-id: o014.aljabr2.chapter7.coalgebra-to-hopf-b.p003
For example, $\munit$ is a Hopf algebra with antipode
$\identity_{\munit}$.

% segment-id: o014.aljabr2.chapter7.coalgebra-to-hopf-b.p004
Interpreted using the convolution in
Remark~\ref{rem:coalgebra-convolution}, an antipode $S$ is precisely a
two-sided inverse of $\identity_A\in\End(A)$ with respect to $\star$. The
antipode in the definition of a Hopf algebra, if it exists, is unique. This
follows immediately by applying the next result to $f=\identity_A$; the
result also shows that morphisms always preserve antipodes.

% segment-id: o014.aljabr2.chapter7.coalgebra-to-hopf-b.pr002
\begin{proposition}\label{prop:antipode-homo}
	Let $A$ and $A'$ be Hopf algebras over $\mathcal{V}$ with antipodes $S$
	and $S'$, respectively. For every Hopf algebra morphism $f:A\to A'$, one
	has $S'f=fS$.
\end{proposition}

% segment-id: o014.aljabr2.chapter7.coalgebra-to-hopf-b.pf002
\begin{proof}
	By the functoriality of convolution $\star$ in
	\eqref{eqn:convolution-functorial}, both maps
% segment-id: o014.aljabr2.chapter7.coalgebra-to-hopf-b.g004
	\[\begin{tikzcd}
		\Hom(A', A') \arrow[r, "{g \mapsto gf}"] & \Hom(A, A')
	\end{tikzcd}, \quad \begin{tikzcd}
		\Hom(A, A) \arrow[r, "{h \mapsto fh}"] & \Hom(A, A')
	\end{tikzcd}\]
	are homomorphisms of convolution monoids. Thus $S'f$ and $fS$ are both
	convolution inverses of $f\in\Hom(A,A')$.
\end{proof}

% segment-id: o014.aljabr2.chapter7.coalgebra-to-hopf-b.p005
Because the braiding on $\mathcal{V}$ gives a braiding on
$\mathcal{V}^{\opp}$, the concept of a bialgebra is plainly self-dual. The
next result shows that the definition of an antipode is self-dual as well.

% segment-id: o014.aljabr2.chapter7.coalgebra-to-hopf-b.pr003
\begin{proposition}\label{prop:bialgebra-duality}
	If $(A,\mu,\eta,\Delta,\epsilon)$ is a bialgebra over $\mathcal{V}$, then
	$(A,\eta,\mu,\epsilon,\Delta)$ is a bialgebra over
	$\mathcal{V}^{\opp}$. If the Hopf algebra
	$(A,\mu,\eta,\Delta,\epsilon)$ has antipode $S$, then $S^{-1}$ is an
	antipode of $(A,\eta,\mu,\epsilon,\Delta)$ as a bialgebra over
	$\mathcal{V}^{\opp}$; in particular, the latter object is a Hopf algebra
	over $\mathcal{V}^{\opp}$.
\end{proposition}

% segment-id: o014.aljabr2.chapter7.coalgebra-to-hopf-b.pf003
\begin{proof}
	The diagram in Definition~\ref{def:Hopf-algebra} is self-dual, but the
	direction of $S$ is reversed.
\end{proof}

% segment-id: o014.aljabr2.chapter7.coalgebra-to-hopf-b.pr004
\begin{proposition}\label{prop:Hopf-antiautomorphism}
% segment-id: o014.aljabr2.chapter7.coalgebra-to-hopf-b.x003
	\index[sym1]{Acop@$A_{\mathrm{cop}}, A^{\mathrm{op}}_{\mathrm{cop}}$}
	For a bialgebra $A$, in accordance with
	Definition~\ref{def:opposite-braided-algebra},
% segment-id: o014.aljabr2.chapter7.coalgebra-to-hopf-b.l002
	\begin{compactitem}
% segment-id: o014.aljabr2.chapter7.coalgebra-to-hopf-b.i003
		\item reversing the multiplication of $A$ using the braiding $c(A,A)$
		produces an algebra $A^{\mathrm{op}}$;
% segment-id: o014.aljabr2.chapter7.coalgebra-to-hopf-b.i004
		\item reversing the comultiplication of $A$ using $c(A,A)^{-1}$
		produces a coalgebra $A_{\mathrm{cop}}$.
	\end{compactitem}
	Together, these two operations produce a bialgebra
	$A^{\mathrm{op}}_{\mathrm{cop}}$. If $A$ has an antipode $S$, then $S$ is
	also an antipode of $A^{\mathrm{op}}_{\mathrm{cop}}$, and in this case
	$S:A\rightiso A^{\mathrm{op}}_{\mathrm{cop}}$ is an isomorphism of
	bialgebras.
\end{proposition}

% segment-id: o014.aljabr2.chapter7.coalgebra-to-hopf-b.pf004
\begin{proof}
	A routine verification shows that $A^{\mathrm{op}}_{\mathrm{cop}}$ is
	indeed a bialgebra with antipode $S$; we omit the details. The remaining
	assertion has two parts. First, $S$ preserves the unit and counit. This is
	easy: in Proposition~\ref{prop:antipode-homo}, take $A=\munit$ and
	$A'=\munit$, respectively.

	Second, $S$ preserves the multiplication and comultiplication. The proof
	in general is more involved; see \cite[\S 9, Proposition 2]{JS91} or
	\cite[Proposition 1.22]{AM10}. Here we only sketch the special case
	$\mathcal{V}=\Bbbk\dcate{Mod}$, $\otimes=\otimes_{\Bbbk}$ with its
	standard braiding. This is the only case that will be needed later.

	Recall from \S\ref{sec:alg-in-monoidal-cat} that $A\otimes A$ is naturally
	a coalgebra. Equip $\Hom(A\otimes A,A)$ with the convolution from
	Remark~\ref{rem:coalgebra-convolution}, denoted by $\ast$, and consider the
	elements
% segment-id: o014.aljabr2.chapter7.coalgebra-to-hopf-b.q002
	\[ f := \mu, \quad g := \mu c(A, A) (S \otimes S), \quad h := S \mu. 
	\]
	It is enough to prove that
	$h\ast f=\eta_A\epsilon_{A\otimes A}=f\ast g$. This equality will show
	that $f$ is invertible in the monoid
	$(\Hom(A\otimes A,A),\ast)$, and hence that $h=g$. Let $a,b\in A$ and
	write the expansions
% segment-id: o014.aljabr2.chapter7.coalgebra-to-hopf-b.g005
	\begin{gather*}
		\Delta(a) = \sum_i a_i^{(1)} \otimes a_i^{(2)}, \quad
		\Delta(b) = \sum_j b_j^{(1)} \otimes b_j^{(2)}.
	\end{gather*}
	Write $\mu$ simply as multiplication. From the standard braiding on
	$\Bbbk\dcate{Mod}$, one obtains
% segment-id: o014.aljabr2.chapter7.coalgebra-to-hopf-b.g006
	\begin{align*}
		(h \ast f)(a \otimes b) & = \sum_{i, j} h\left(a^{(1)}_i \otimes b^{(1)}_j\right) f\left(a^{(2)}_i \otimes b^{(2)}_j\right) \\
		& = \sum_{i, j} S\left( a^{(1)}_i b^{(1)}_j \right) a^{(2)}_i b^{(2)}_j \\
		& = (\underbracket{S \star \identity_A}_{\text{convolution on }\End(A)})(ab) = \eta_A \epsilon_A (ab), \\
		(f \ast g)(a \otimes b) & = \cdots = \sum_{i, j} a^{(1)}_i b^{(1)}_j S\left( b^{(2)}_j \right) S\left( a^{(2)}_i \right) \\
		& = \sum_i a^{(1)}_i \eta_A \epsilon_A(b) S\left( a^{(2)}_i \right) \\
		& = \eta_A \epsilon_A (a) \cdot \eta_A \epsilon_A (b) = \eta_A \epsilon_A(ab), \\
		\eta_A \epsilon_{A \otimes A}(a \otimes b) & = \eta_A\left( \epsilon_A(a)\epsilon_A(b) \right) = \eta_A \epsilon_A(ab),
	\end{align*}
	as claimed. Thus $S$ preserves the multiplication of the bialgebra. A
	similar argument shows that $S$ preserves the comultiplication.
\end{proof}

% segment-id: o014.aljabr2.chapter7.coalgebra-to-hopf-b.p006
In particular, if $\mathcal{V}$ is a symmetric monoidal category, then $S^2$
is an automorphism of $A$.

% segment-id: o014.aljabr2.chapter7.coalgebra-to-hopf-b.c001
\begin{corollary}\label{prop:Hopf-involution}
	Let the Hopf algebra $A$ have antipode $S$. If $A$ is commutative or
	cocommutative, then $S^2=\identity$.
\end{corollary}

% segment-id: o014.aljabr2.chapter7.coalgebra-to-hopf-b.pf005
\begin{proof}
	With respect to the convolution $\star$ on $\End(A)$
	(Remark~\ref{rem:coalgebra-convolution}), we have
% segment-id: o014.aljabr2.chapter7.coalgebra-to-hopf-b.q003
	\[ S \star S^2 = \mu (S \otimes S^2) \Delta = \mu (S \otimes S)(\identity \otimes S) \Delta. \]

	If $A$ is commutative, then $\mu c(A,A)=\mu$, and
	Proposition~\ref{prop:Hopf-antiautomorphism} transforms the right-hand side
	above into
% segment-id: o014.aljabr2.chapter7.coalgebra-to-hopf-b.q004
	\[ \mu c(A, A) (S \otimes S) (\identity \otimes S)\Delta = S\mu (\identity \otimes S)\Delta = S(\identity \star S). \]
	The definition of an antipode says that
	$\identity\star S=\eta\epsilon$, while the same proposition gives
	$S\eta=\eta$; the result is $\eta\epsilon$.

	If $A$ is cocommutative, then $c(A,A)\Delta=\Delta$, and functoriality of
	the braiding transforms $S\star S^2$ into
% segment-id: o014.aljabr2.chapter7.coalgebra-to-hopf-b.q005
	\[ \mu (S \otimes S)(\identity \otimes S) c(A, A) \Delta = \mu c(A, A) (S \otimes S)(S \otimes \identity)\Delta, \]
	which the same argument then transforms into
	$S(S\star\identity)=\eta\epsilon$.

	Thus $S^2$ is a convolution inverse of $S$, so $S^2=\identity$.
\end{proof}

% segment-id: o014.aljabr2.chapter7.coalgebra-to-hopf-b.p007
All the following examples take $\mathcal{V}=\Bbbk\dcate{Mod}$ and
$\otimes=\otimes_{\Bbbk}$, where $\Bbbk$ is a commutative ring.

% segment-id: o014.aljabr2.chapter7.coalgebra-to-hopf-b.ex001
\begin{example}\label{eg:group-Hopf-algebra}
	Let $G$ be a group. Form the group $\Bbbk$-algebra $\Bbbk[G]$.
	Example~\ref{eg:monoid-bialgebra} gives $\Bbbk[G]$ a bialgebra structure.
	Define the $\Bbbk$-module automorphism $S:\Bbbk[G]\to\Bbbk[G]$ by
% segment-id: o014.aljabr2.chapter7.coalgebra-to-hopf-b.q006
	\[ S(g) = g^{-1}, \quad g \in G. \]
	It is easy to verify that $S$ satisfies the commutative diagram required
	of an antipode. Thus $\Bbbk[G]$ is a Hopf algebra.
\end{example}

% segment-id: o014.aljabr2.chapter7.coalgebra-to-hopf-b.ex002
\begin{example}[H.\ Hopf]
% segment-id: o014.aljabr2.chapter7.coalgebra-to-hopf-b.x004
	\index{H-space, H-group}
	Write $\{\mathrm{pt}\}$ for a one-point set. Let a topological space $X$
	be equipped with continuous maps $m:X\times X\to X$ (multiplication) and
	$e:\{\mathrm{pt}\}\to X$ (equivalently, a choice of an element of $X$,
	the unit, also denoted by $e$), such that the two maps satisfy, up to
	homotopy, the properties required for associativity of multiplication and
	for a unit. Then $(X,m,e)$ is called an \emph{H-space}. If a continuous map
	$i:X\to X$ (taking inverses) is also chosen and satisfies, up to homotopy,
	the properties required of inverses, then $(X,m,e,i)$ is called an
	\emph{H-group}.

	A topological monoid (or topological group) is of course an H-space (or
	H-group), but the H-space and H-group conditions are much weaker. A natural
	example is a loop space. For a topological space $M$ with a chosen base
	point $x_0\in M$, a continuous map $\gamma:[0,1]\to M$ satisfying
	$\gamma(0)=x_0=\gamma(1)$ is called a \emph{loop} in $M$. All loops
	naturally form a topological space $\Omega(M,x_0)$. If $m$ is taken to be
	end-to-start concatenation of two loops, $e$ the constant map at $x_0$,
	and $i$ reversal of the direction of a loop, then intuitively
	$\Omega(M,x_0)$ forms an H-group. The key point is that associativity and
	the inverse properties hold only up to homotopy.

	Next, consider the full subcategory $\cate{S}$ of $\cate{Top}$ consisting
	of all CW complexes\footnote{This is a topological term, somewhat different
	from the complexes defined in linear algebra.}; it contains spaces commonly
	encountered in geometry, such as manifolds. Let $\Bbbk$ be a field. The
	cohomology functor with coefficients in $\Bbbk$, namely $\Hm^\bullet$, is a
	functor from $\cate{S}^{\opp}$ to the category of graded $\Bbbk$-vector
	spaces $\cate{Vect}(\Bbbk)^{\Z_{\geq0}}$. Equip
	$\cate{Vect}(\Bbbk)^{\Z_{\geq0}}$ with the braiding determined by the
	Koszul sign rule \eqref{eqn:Koszul-braiding} (in the notation there, take
	$\epsilon(a)=a\bmod\;2$). The cup product $\mu:=\cup$ on cohomology makes
	the functor $\Hm^\bullet$ factor through
	$\cate{CAlg}\left(\cate{Vect}(\Bbbk)^{\Z_{\geq0}}\right)$; in short,
	cohomology is a commutative algebra over
	$\cate{Vect}(\Bbbk)^{\Z_{\geq0}}$.

	On the other hand, give $\cate{S}$ the monoidal structure determined by
	the product of spaces, with $\{\mathrm{pt}\}$ as unit. The Künneth formula
	in topology shows that $\Hm^\bullet$ is a monoidal functor. Moreover,
	homotopic maps induce the same map on cohomology. Consequently, if
	$X\in\Obj(\cate{S})$ is an H-space, then $\Hm^\bullet(X)$ is also a
	coalgebra over $\cate{Vect}(\Bbbk)^{\Z_{\geq0}}$.

	Readers familiar with algebraic topology will readily conclude that the
	algebra and coalgebra structures on $\Hm^\bullet(X)$ are compatible and
	hence give a bialgebra over
	$\cate{Vect}(\Bbbk)^{\Z_{\geq0}}$. If, in addition, $X$ is an H-group and
	the automorphism induced by the map $i$ is denoted by
	$S\in\Aut(\Hm^\bullet(X))$, then $\Hm^\bullet(X)$ becomes a Hopf algebra.

	In fact, the multiplication (cup product) on $\Hm^\bullet(X)$ is precisely
	the image of the diagonal embedding $\mathrm{diag}:X\to X\times X$ under
	the Künneth formula, while its unit is the image of
	$X\to\{\mathrm{pt}\}$. Thus all the properties required of the bialgebra
	and antipode can be verified at the level of topological spaces, up to
	homotopy. For example, the antipode identity
	$\mu(S\otimes\identity)\Delta=\eta\epsilon$ comes from
% segment-id: o014.aljabr2.chapter7.coalgebra-to-hopf-b.q007
	\[ \left[ X \xrightarrow{\mathrm{diag}} X \times X \xrightarrow{(i, \identity)} X \times X \xrightarrow{m} X\right] \;\text{being homotopy equivalent to}\; \left[ X \to \{\mathrm{pt}\} \xrightarrow{e} X \right]. \]

	Thus the cohomology of an H-group is naturally a commutative Hopf algebra
	with a rich structure. This was Hopf's original motivation for studying
	such algebras.
\end{example}

% segment-id: o014.aljabr2.chapter7.coalgebra-to-hopf-b.p008
The exercises will provide more examples of Hopf algebras. Most are
commutative or cocommutative, but there are also many Hopf algebras having
neither property. The most important class among them is that of
\emph{quantum groups}. The relevant discussion requires a foundation in Lie
theory and is not treated in this book.
